\documentclass[11pt]{article}
\usepackage[margin = 2cm]{geometry}
\usepackage{eulervm}

\usepackage{hyperref}
\usepackage{xcolor}
\usepackage{amsmath}
\usepackage{amsthm}
\usepackage{amssymb}
\usepackage{physics}
\usepackage{dsfont}

\newtheorem{defn}{Def}
\newtheorem{prop}{Prop}
\newtheorem{lem}{Lemma}
\newtheorem{cor}{Cor}
\newtheorem{thm}{Theorem}

\theoremstyle{remark}
\newtheorem*{rmk}{Remark}

\newcommand{\catname}[1]{{\normalfont\textbf{#1}}}
\newcommand{\e}{\mathrm{e}}

\begin{document}
\noindent\huge{The maximum entropy principle} \\
\normalsize\textbf{Blog Post \today} \\

Something that fascinated me in undergrad stat mech was that the Shannon entropy
falls out of the canonical ensemble upon some relatively reasonable assumptions.
Here, we will go both ways; First, we will explore the natural properties of the
Shannon entropy from the perspective of \cite{baez}, and then use the method of Lagrange multipliers to derive the canonical
ensemble from the maximum entropy principle. Then, we will go the other way
following \cite{reif}, and
show that the Shannon (or Von Neumann) entropy comes out of the canonical ensemble.

\section{Shannon entropy}
The Shannon entropy is natural in the sense that it satisfies the following
properties, formalized by Faddeev \cite{faddeev}:
\begin{thm}
  Suppose that $H: p \to \mathbb R$, where $p$ is a discrete probability measure and $H$
  satisfies the following axioms:
\begin{itemize}
\item Permutation invariance: $H(p) = H({}^{\sigma}p)$ for any $\sigma \in S_n$ denoting a permutation
  ${}^{\sigma}p(i) = p(\sigma(i))$
\item $H(p, 1-p)$ is a continuous function of $p$.
\item $H(\ast) = 0$, where $\ast$ is the measure space consisting of a single element.
\item Recursivity: $H(\lambda p_1, (1-\lambda)p_1, p_2, \dots, p_n) = p_1 H(\lambda, 1-\lambda) + H(p)$.
\end{itemize}
Then $H = -c \sum_{i}p(i)\log(p(i))$, where $c$ is some positive constant.
\end{thm}

The first condition is eminently reasonable; Since $H$ is a measure of
information assigned to a distribution, permuting the outcomes should have no effect.
The continuity condition says that changing the distribution by a small amount should
not greatly affect the entropy. The third condition holds that there is no
information in a non-random outcome.

The last condition is arguably the most important; define $\phi(p) = H(\{p\})$.
By iterating condition 3, we may find that $\phi(pq) = \phi(p) + \phi(q)$. It is
tedious to show due to the fact that $p,q$ are inverses of integers, but the
unique $\phi$ satisfying this requirement is $\phi(p) \propto \log(p)$. This
last condition, known as the grouping property or recursivity, seems more difficult to interpret than the previous ones---we
will follow \cite{baez}, which shows that the entropy difference satisfies nice
properties and illuminates this third requirement.

Define \catname{FinProb} to be the category of probability measures on finite
spaces. The morphisms of \catname{FinProb} are defined to be measure-preserving
maps. A map $T: (X, p) \to (Y, q)$ is said to be measure-preserving if $T_\ast
p(A) = q(T^{-1}(A)) = p(A)$, where $A \subseteq X$.
Importantly, given $(X, p), (Y,q)$ and $\lambda \in [0,1]$, we can construct a probability measure $\lambda p
\oplus (1-\lambda)q$ on $X \sqcup Y$, which intuitively represents sampling from
$p$ with probability $\lambda$ and from $q$ with probability $1-\lambda$. If $T_1 : (X_1, p_1) \to (X_2, p_2)$ and
$T_2:(Y_1, q_1) \to (Y_2, q_2)$ are measure-preserving maps, then $T_1 \oplus T_2$ is a measure-preserving map $(X_1\sqcup X_2, \lambda p_1
\oplus (1-\lambda)p_2) \to (Y_1 \sqcup Y_2, \lambda q_1 + (1-\lambda)q_2)$. 
Instead of $H$, we will define a functor $F:\catname{FinProb} \to \mathbb [0, \infty)$
called the \emph{entropy loss}:

\begin{defn}
  Let $F: \catname{FinProb} \to \mathbb [0, \infty]$ satisfy the following axioms:
  \begin{itemize}
  \item Functoriality: $F(T_1 \circ T_2) = F(T_1) + F(T_2)$
    \item Convex-linear: $F(T_1 \oplus T_2) = \lambda F(T_1) +
      (1-\lambda)F(T_2)$ if $T_1 \oplus T_2$ is a map on $\lambda p \oplus (1-\lambda)q$
     \item Continuity: We say that a sequence
     $T_n : (X_n, p_n) \to (Y_n, q_n)$ converges to $T: (X, p) \to (Y, q)$ if
     $T_n = T$, $X_n= X$, and $Y_n = Y$ for large enough $n$ and $p_n,q_n \to p,
     q$ in the pointwise sense, and $F$ is continuous if $F(T_n) \to F(T)$.

     \end{itemize}
     \end{defn}

The interpretation of the first condition is obvious; if we perform processes
$T_1$ and $T_2$ in sequence, we lose $F(T_1)$ of information at the first step
and $F(T_2)$ information at the second step, so we want to say that $F(T_1) +
F(T_2)$ information is lost overall. The second condition amounts to saying that
if we perform process $T_1$ with probability $\lambda$ and $T_2$ with probability
$1-\lambda$, then our information loss is the average loss between both
processes. Our main goal will be to show that these conditions are equivalent to
the recursivity condition of Faddeev. It was recognized by Renyi that this averaging
condition is stronger than one might like: Replacing this condition with a more
general one leads to the Renyi entropies \cite{renyi}.

\begin{thm}
  If $(X,p), (Y,q)$ are discrete measure spaces and $T: X \to Y$ is a
  measure-preserving map, then $F(T) = c(H(p) - H(q))$, where $H$ is the Shannon
  entropy and $c > 0$ is a constant.
\end{thm}

\begin{proof}
First, define the ``crunch map'' $!_q: (X, q) \to \{\ast\}$ via
$!_q(x) = \ast$. Put $H(q) = F(!_q)$ to be the entropy of $q$, and we will show
that $H$ satisfies the hypotheses of Faddeev's theorem. First, for any $T: (X,
p) \to (Y,q)$, we have $!_p
= !_q \circ T$, so $H(p) = F(!_p) = F(!_q) + F(T) = H(q) + F(T)$. This shows
that $F(T) = H(p) - H(q)$.

Lastly, we can show why the convex-linearity
condition is equivalent to the additivity condition. Given spaces $(X_1, p_1), \dots,
(X_n, p_n)$ and a measure $p$ on $[n]$, we obtain a measure on $\hat X = \sqcup_{i=1}^n
X_i$ via $\hat p(x_i) = p(i)p_i(x_i)$. We can then form the map $\pi : \hat X
\to [n]$ via $x_i \mapsto i$. Convex-linearity then requires that
\begin{align}
F(\pi) = F(\bigoplus_{i} !_{p_i}) = \sum_i p(i)F(!_{p_i}) = \sum_i p(i)H(p_i)
\end{align}
But we also have $F(f) = H(\hat p) - H(p)$. This implies condition 3.
\end{proof}
What we have proven is in fact stronger than recursivity, and it
is frequently referred to as strong additivity. \\

For a physicist, it is nice to see how one might stumble upon the Shannon
entropy without a lengthy derivation. The quick and dirty answer is this; If all
messages are sent with probability $p$, then the message contains $\log_2(1/p)$
bits. Therefore if one receives a message with probability $p$, one learns
$\log_2(1/p)$ bits of information. This quantity is known as the `surprisal',
and the entropy is just the average amount of surprisal. 

\section{Canonical ensemble from Lagrange multipliers}
Consider a statistical ensemble taking on any one of a number of discrete states
$r$ with probability $p_r$ and energy $E_r$. Let $S = -\sum_{r}p_r\ln(p_r)$ be
the entropy. Suppose that we know the average energy $\overline E = \sum_r E_r
p_r$. We assume that in equilibrium, the system is equally likely to occupy any
state consistent with the average energy, which equates to positing that $S$ is
maximal with respect to the constraint $\sum_r E_r p_r = \overline E$. We
additionally have the constraint $\sum_r p_r = 1$, so we
introduce two Lagrange multipliers $\beta, \gamma$:
\begin{align}
  \pdv{p_r}S + \beta \pdv{p_r}\overline E + \gamma \pdv{p_r}\sum_{r}p_r &= 0 \\
  \ln(p_r) + 1 - \beta E_r + \gamma = 0 \\
  p_r & \frac{\e^{-\beta E_r}}{Z}
\end{align}
Where $Z = \sum_r \e^{-\beta E_r}$ is the partition function. Similarly, if the number $\overline N = \sum_r N_r p_r$ is allowed to vary as
well, we have
\begin{align}
  \pdv{p_r}S + \beta \pdv{p_r}\overline E + \alpha \pdv{p_r}\overline N + \gamma \pdv{p_r}\sum_r p_r &= 0 \\
  \ln(p_r) + 1 + \beta E_r + \alpha N_r+ \gamma = 0 \\
  p_r = \frac{\e^{-\beta E_r - \alpha N_r}}{\mathcal Z}
\end{align}
which gives us the grand canonical ensemble. Sticking with the canonical
ensemble for simplicity, we have
\begin{align}
  S &= -\sum_r p_r \ln p_r = \sum_r p_r (\beta E_r + \ln Z) = \beta \overline E + \ln Z
\end{align}
From this, we can
derive all of the typical relationships of statistical mechanics. Of particular
interest is the following:
\begin{align}
  \dd S &= \dd (\beta \overline E + \ln Z)\\
  &= \dd \beta \overline E + \beta \dd \overline E + \pdv{\ln Z}{x}\Big|_\beta\dd x + \pdv{\ln
    Z}{\beta}\Big|_x\dd \beta
\end{align}
This is from the fact that $E$ is a function of some external parameters $x$
(like volume or trap depth, for instance), and thus $Z$ is a function of $x$ and
$\beta$.
Then we notice that
\begin{align}
\pdv{\ln Z}{\beta} = \frac{1}{Z}\pdv{Z}{\beta} = \frac{-\beta}{Z}\sum_r E_r \e^{-\beta E_r} = -\beta \overline E
\end{align}
Next, we have
\begin{align}
\pdv{\ln Z}{x}\Big|_{\beta}\dd x = \frac{1}{Z}\pdv{Z}{x}\Big|_\beta\dd x = \frac{-\beta}{Z}\sum_r E_r \pdv{E_r}{x}\Big|_\beta\e^{-\beta E_r} \dd x = -\beta \dd \overline W
\end{align}
where we identify $\pdv{E_r}{x}\Big|_{\beta}\dd x = \dd W_r$ as the work done by
the system.
Plugging this in, we now have
\begin{align}
\dd S = \dd \beta \overline E + \beta \dd \overline E - \beta\dd \overline W - \beta \overline E \dd \beta = \beta(\dd \overline E - \dd \overline W) = \beta \dd Q = \frac{\dd Q}{T}
\end{align}
where we defined $\dd Q = \dd \overline E - \dd \overline W$ to be the change in
energy of the system \emph{not} caused by changing external parameters, i.e.,
the change in heat. Thus, we have $\dd \overline E = \dd Q + \dd W = T\dd S - p
\dd V$. This is enough to give us all our normal thermodynamic equations, e.g.
$\pdv{\overline E}{S}\Big|_V = T$.

\section{Canonical ensemble from interaction with a heat bath}
In the above derivation, no reference was made to the interaction with a heat
bath or the number of accessible microstates. To furnish this picture, assume a
system $A$ is interacting with a much larger heat bath $A'$. The maximum entropy
principle is equivalent to requiring that the probability $p_r$ of a microstate
$r$ for the system $A$ is equivalent to $\frac{\Omega'(E_0-E)}{\Omega_0(E_0)}$,
where $E_0 = E + E'$ is the total energy of the system, and $\Omega_0, \Omega,
\Omega'$ are the number of microstates available to each system at a given
energy. The entropy is maximized by a uniform distribution. If system $A$ is in
the state $r$ with energy $E_r$, then the \emph{bath} is equally likely to
occupy any state
$r'$ with $E_{r'}' = E_0-E$ Since $E \ll E_0$ when $A$ is much smaller than
$A'$, we have
\begin{align}
  p_r &= \frac{\Omega'(E_0-E_r)}{\Omega_0(E_0)} \\
  \ln p_r &= C + \ln(\Omega' (E_0-E_r)) \approx C + \ln(\Omega'(E_0)) - \pdv{\Omega'(E_0)}{E}E_r \equiv C' - \beta' E_r \\
  p_r &\propto \e^{-\beta' E_r}
\end{align}
Next, we again appeal to the maximum entropy principle: We have $\Omega_0(E_0) =
\Omega'(E_0-E)\Omega(E)$, and so $\ln \Omega(E_0) = \ln \Omega'(E_0-E) + \ln
\Omega'(E)$. Since $E$ will be such that $\Omega_0(E_0)$ is maximized, we will have
\begin{align}
\pdv{E} \ln \Omega_0(E_0) = \pdv{\ln \Omega(E)}{E} - \pdv{\ln \Omega'(E_0 -E)}{E} = \beta - \beta' = 0
\end{align}
Now we can identify $p_r \propto \e^{-\beta E_r}$, the canonical ensemble. In particular,
with volume fixed, we have $\dd E = \dd Q$, and 
\begin{align}
\dd\ln \Omega(E) = \ln \Omega(E+\dd Q) - \ln \Omega(E) \approx \pdv{\ln \Omega(E)}{E}\dd Q = \beta \dd Q
\end{align}
From before, we have $\beta \dd Q = \dd S$, so $S = \ln \Omega(E)$. All the steps we
did previously are reversible; Assuming that $S = \ln \Omega(E)$, we then find
$\dd S = \beta \dd Q$. It follows that $S = \beta \overline E + \ln Z$,
and then from there we find $S = \sum_{r}p_r \ln p_r$.

\bibliography{refs}
\bibliographystyle{plain}

\end{document}


